\sectionguard
\section*{The Propers: Notable and Quotable}

Four verified afterlives of wording from the scriptural propers, each a use
that moves the phrase into a register its author did not occupy. Straight
exegesis, devotional reuse, and bare quotation are excluded; they belong in the
commentary above. Two of the four come from one Aramaic word, and that
concentration is itself the finding: of everything on these three missal pages,
\latin{Ephphetha} is what the world took and put to other work. The audit
behind the gallery --- including a fifth verified candidate held back, and the
searches that came up empty --- is in \texttt{research/scope.md}.

\begin{afterlife}{Epistle, 1 Cor. 15:8 --- William Morris makes an apostle's apology into an aesthete's manifesto, 1868}
\textbf{Appointed phrase.} \latin{Novíssime autem ómnium tamquam abortívo, visus
est et mihi}; Douay--Rheims, ``And last of all, he was seen also by me, as by
one born out of due time.''

\textbf{Later use and locus.} William Morris, \work{The Earthly Paradise}
(London: F.~S. Ellis, 1868), the prefatory ``Apology,'' stanza 4, on p.~2 of the
first edition: ``Dreamer of dreams, born out of my due time, \textbar\ Why
should I strive to set the crooked straight?'' Read in the archive.org scan of
the first edition.

\textbf{The turn.} Paul's phrase is an apology for the violence and lateness of
his own apostolic vision, made in the act of claiming apostolic standing.
Morris keeps the idiom, inserts one word --- \textit{my} --- and inverts the
claim. Untimeliness now means being a medieval dreamer stranded in the
industrial nineteenth century, and where Paul's untimely birth is the ground of
a commission, Morris's is the ground of a disclaimer: he explicitly denies any
power to set the crooked straight. The Apostle's least of the apostles becomes
the idle singer of an empty day.

\textbf{Rights and limit.} The first edition (1868) and Morris (d.~1896) are
long out of copyright. The dependence is on the naturalised English idiom
``born out of due time,'' shared by the Douay and the King James, not on this
Mass lection; Morris cites nothing. The ``Apology'' is absent from the Project
Gutenberg Part~II e-text, which is why the first-edition scan was used.
\end{afterlife}

\begin{afterlife}{Gospel, Mk. 7:34 --- ``Effeta, The Body Deodorant,'' New Jersey, 1912}
\textbf{Appointed phrase.} \latin{Ephphetha, quod est adaperíre}, and the
baptismal formula that quotes it: \latin{Ephpheta, quod est, Adaperire \dots\ In
odorem suavitatis} (\work{Rituale Romanum}, Pustet 1872, pp.~11--12).

\textbf{Later use and locus.} A newspaper advertisement for the Willis Pike Co.\
Laboratories, 129 W.\ 31st St., New York, in the \work{Perth Amboy Evening
News} (N.J.), 19 June 1912, ed.~2, p.~2: ``For That Embarrassing Odor of
Perspiration Get a Box of Effeta, The Body Deodorant. `Be thou opened unto odor
of sweetness.''' The campaign runs in the same paper from 29 May to 22 June
1912 and reappears in the \work{Norwich Bulletin} (Conn.), 22 June 1912, p.~6.
Read in the Library of Congress \work{Chronicling America} page text.

\textbf{The turn.} The product name is the Latin form of the Gospel word ---
and, at that, the form Bede's lemma uses --- and the slogan splices Christ's
command to the deaf man onto the rite's own \latin{in odorem suavitatis}, the
words said over the nostrils. What is opened in Mark is a pair of ears and a
tongue; what is opened in the advertisement, for twenty-five cents a box, is
the pores. The copywriter is not being blasphemous so much as commercially
literal: he has found a sacramental formula about opening and odour and sold it
as a toiletry.

\textbf{Rights and limit.} United States newspapers of 1912 are public domain.
The copywriter's proximate source is the baptismal rite rather than the Mass
lection --- but the rite is quoting this Sunday's Gospel, and St.~Ambrose says
so. The 19 June issue is the one in which the scriptural slogan line is legible
in the page text, and the wording was read there in context.
\end{afterlife}

\begin{afterlife}{Introit psalm verse, Ps. 67:2 --- Cromwell at Dunbar, 3 September 1650}
\textbf{Appointed phrase.} \latin{Exsúrgat Deus, et dissipéntur inimíci eius};
Douay--Rheims, ``Let God arise, and let his enemies be scattered.'' It is the
verse the missal appoints as the Introit's psalm verse.

\textbf{Later use and locus.} At the moment the Scottish line broke, with the
sun coming up over the sea, Cromwell was heard to exclaim ``Now let God arise,
and his enemies shall be scattered.'' The eyewitness is Captain John Hodgson,
\work{Original Memoirs \dots\ Life of Sir Henry Slingsby, and Memoirs of Capt.\
Hodgson} (Edinburgh, 1806), pp.~147--48, whose editor footnotes the reference
as ``Psalm lxviii.\ v.~1'' and who adds, a few lines later, ``I heard him say,
`I profess they run!''' Relayed with the psalter version by Thomas Carlyle,
\work{Oliver Cromwell's Letters and Speeches} (Tauchnitz, 1861), vol.~2,
pp.~332--33. Three years later, opening Barebone's Parliament on 4 July 1653,
Cromwell read the whole psalm aloud; Carlyle notes ``We remember it ever since
Dunbar morning'' (vol.~3, pp.~171--72).

\textbf{The turn.} In the Mass the verse is sung by a congregation as the
priest comes to the altar. At Dunbar it is not a prayer at all but a
battlefield exclamation --- the shout of a commander who has just watched a
cavalry charge succeed, using the psalm to certify the outcome rather than to
ask for it. Three years later the same psalm is deployed politically, as
scriptural warrant for a revolutionary assembly. The verse travels from
petition to victory-cry to constitutional argument.

\textbf{Rights and limit.} Both sources are public domain. Cromwell's wording is
the English psalter and King James tradition, not the Douay. Hodgson's memoir is
remembered speech first printed in 1806, not a contemporary dispatch; Cromwell's
own letter about Dunbar does not contain the phrase. The 1806 page text garbles
the word ``scattered,'' and the reading here is controlled by the editor's
footnote and by Carlyle's transcription.
\end{afterlife}

\begin{afterlife}{Gospel, Mk. 7:34 --- a command to the deaf becomes the name deaf institutions gave themselves}
\textbf{Appointed phrase.} \latin{Ephphetha}, ``Be thou opened.''

\textbf{Later use and locus.} (a) The Ephphatha Guild, a body of ``between
fifty and sixty'' deaf-mutes attached to St.~Luke's, Scranton, with an
affiliated secular ``Penny Club'' and its own fundraising fairs ---
\work{The Scranton Tribune}, 14 November 1895, p.~3, under the heading
``EPHPHATHA GUILD SERVICES,'' which supplies the derivation itself: the guild
``takes its name from the Greek word `Ephphatha' `Be opened' used by the Savior
in healing a deaf mute.'' (b) The Ephphatha Missions for the Deaf and Blind and
the Ephphatha Services Board of the American Lutheran Church, with an annual
professional Ephphatha Conference --- Council of Organizations Serving the
Deaf, \work{Council Membership Directory 1969} (ERIC ED032704).

\textbf{The turn.} A word spoken once, to one man, becomes a durable
institutional surname across denominations and three-quarters of a century:
dues, guild fairs, boards, a conference in the calendar. The register shifts
from miracle to administration. And there is a sharper edge: in the decades of
the oralism-versus-sign controversy, deaf people organised themselves under a
word that commands the deaf to hear and the mute to speak, and used it as a
name for their own society rather than as a prescription addressed to them.

\textbf{Rights and limit.} The 1895 newspaper is public domain; the 1969
directory is a freely distributed ERIC document, cited without extended
reproduction. Both bodies are church-affiliated, so they qualify as proper
names of functioning institutions rather than as sermon reuse. The Scranton
paper miscalls the Aramaic word ``Greek,'' which is reproduced above as
printed.
\end{afterlife}
